\section{Extended Attack-Resilience Results}
\label{app:attacks}

Table~\ref{tab:attack-full} reports the ten-scenario single-actor attack
suite, rerun for this paper on the seven-validator testbed ($n{=}7$) with
three repeats per scenario. The repeats of a scenario share one chain, each starting
from \textsc{Anchored}, so they are not independent deployments. A3 and A4 are
the cases of the collusion suite (Section~\ref{subsec:rq1}) with one Byzantine
validator. The column \emph{Byz.\ VP} gives the share of voting power held by
the misbehaving validator: one of seven validators is 14.3\%, and the attacker
containers of A1 and A2 are non-validators (0\%).

\begin{table}[H]
\caption{Single-actor attack suite at $n{=}7$ (live; three repeats each; safety
held and zero AppHash divergence in all 30 repeats)}
\label{tab:attack-full}
\centering
\footnotesize
\begin{tabular}{@{}p{0.04\linewidth}p{0.17\linewidth}p{0.07\linewidth}p{0.55\linewidth}@{}}
\toprule
\textbf{\#} & \textbf{Attack} & \textbf{Byz.\ VP} & \textbf{Attack-specific signal} \\
\midrule
A1 & Eclipse (10 attackers, one subnet) & 0\% & Filter held in 3/3: peak 8 peers in the target subnet, equal to the limit of 8; nil-prevote ratio 0.0. The FSM left \textsc{Anchored} for \textsc{Suspicious} in 2/3 repeats \\
A2 & Multi-subnet Sybil (12 attackers, 4 subnets) & 0\% & Filter held in 3/3 (peak 8 of 8). The FSM left \textsc{Anchored} for \textsc{Suspicious} in 2/3 repeats \\
A3 & Data withholding & 14.3\% & Not rejected by honest validators; replay clean \\
A4 & Forged Bitcoin receipt & 14.3\% & Rejected by honest validators (216 rejections over 3 repeats) \\
A5 & Withdrawal during \textsc{Sovereign} & --- & Blocked in 3/3; the chain did not reach \textsc{Sovereign} and the transaction never committed \\
A6 & Malicious proposer (forged FSM state) & 14.3\% & Safety held in 3/3 (Section~\ref{subsec:rq1}) \\
A7 & Censorship & 14.3\% & Censoring observed in 3/3; liveness held; honest reject signals stayed at 0 \\
A8 & Combined (A6 + A1) & 14.3\% & Safety held under simultaneous overlap \\
A9 & Double-signing & 14.3\% & 6/6 honest validators logged real duplicate-vote evidence in 3/3 repeats, detected one block after the offense \\
A10 & Timeout flood & 14.3\% & Cadence held: 0.754$\to$0.785, 0.791$\to$0.791 and 0.816$\to$0.783 blocks/s (baseline$\to$flood) \\
\bottomrule
\end{tabular}
\end{table}

The \textsc{Suspicious} excursions in A1 and A2 occur in the attack phase and
again in the recovery phase, after the attacker is gone, and the baseline phase
was always \textsc{Anchored}. A control run of A1 without an attacker (three runs, the same phases and timing,
seven validators) shows the same excursions: the validators were in
\textsc{Suspicious} in 17\% to 19\% of the samples and never in \textsc{Sovereign}. The
excursions of A1 are therefore the sensor noise of the testbed, which the collusion
suite's control c0 also shows, and not an effect of the attack. A2 has no control of
its own. A3 and
A4 were instrumented only for safety, divergence and rejections, without an
attack-specific metric such as A5's blocked-withdrawal flag or A9's detection
latency. No scenario drove block production toward zero.

\FloatBarrier

\subsection*{Function-level checks of the weighted aggregation}

The testbed and the TLA+ model give every
validator the same power, where weight and count coincide, so a defect that
confuses them would not appear in any run above. We therefore tested the
aggregation function alone (\texttt{tests/weighted}) on three weight vectors that
sum to 100, $[30,20,15,10,10,10,5]$, $[26,19,17,13,11,8,6]$ and $[34,33,33]$, and
on seven equal weights. Against a reference weighted median (the element at
position $\lfloor T/2\rfloor$ of the multiset in which each value is repeated by
its power) the function agreed on 1{,}500 random quorum subsets with random
values for every vector, and a verified power of exactly $\lfloor 2W/3\rfloor{+}1$
was accepted and one below it rejected for every subset. For every Byzantine
coalition and every set of certificates the proposer may omit while keeping a
quorum, with the adversary reporting the lowest value, the highest, or a split,
the aggregate stayed within the honest values and honest booleans were not
flipped whenever the coalition's power was below $W/3$ (up to 30 of 100 for the
first vector and 33 for the other two). A coalition of 35 (20 and 15) in the first
vector, and of 34 in the other two, escaped the honest range by having the proposer omit
enough honest certificates, so the bound is tight. This checks the
implementation of the weighted median and the quorum, not the integration
with the genesis power assignment, which a live run with unequal power would
add.

\textbf{Proposer influence.} For $n$ from 4 to 10 equal-power validators and 30
random sets of distinct honest values each, we ran the aggregation on every
subset the proposer may choose (all subsets of at least $\lfloor 2n/3 \rfloor + 1$
observations). The smallest and largest aggregate matched the order statistics of
Section~\ref{subsec:health-certificate} in every case.

\subsection*{P2P detector sensitivity}

A P2P detector is useful for safety only if its
signal cannot be abused by a minority. We compare the six-bound
profiler (Section~\ref{subsec:sensors}) with a peer-count-only detector in a
Monte Carlo sensitivity study on synthetic attacks (2{,}000 benign and 2{,}000
attack snapshots plus 500 ramps per cell; Section~\ref{subsec:montecarlo}); the
two most diagnostic attacks were then run live (Section~\ref{subsec:rq1}).

\begin{table}[H]
\caption{P2P detector accuracy: six-bound profiler against peer-count-only
(synthetic sensitivity analysis; 2{,}000 benign and 2{,}000 attack snapshots, 500 ramps per cell)}
\label{tab:p2p-detector}
\centering
\footnotesize
\setlength{\tabcolsep}{3pt}
\begin{tabular}{@{}lrrrr@{}}
\toprule
 & \multicolumn{2}{c}{\textbf{Six-bound profiler}} & \multicolumn{2}{c}{\textbf{Peer-count-only}} \\
\cmidrule(lr){2-3}\cmidrule(lr){4-5}
\textbf{Attack} & \textbf{Delay} & \textbf{Evasion} & \textbf{Delay} & \textbf{Evasion} \\
\midrule
P1 Peer-slot exhaustion & 2.23 $\pm$ 0.87 & 0.17 & never & 1.00 \\
P2 Multi-subnet Sybil & 1.00 $\pm$ 0.00 & 0.00 & never & 1.00 \\
P3 Churn-based rotation & 1.70 $\pm$ 0.61 & 0.12 & never & 1.00 \\
P4 Relay-node latency & 2.71 $\pm$ 1.19 & 0.25 & never & 1.00 \\
\bottomrule
\end{tabular}
\\[2pt]
\footnotesize Delay: mean $\pm$ sample standard deviation of the snapshots to
first detection over the 500 ramps of 10 snapshots each (``never'': no ramp
was detected). Evasion: largest attack intensity on
the 0--1 scale at which the detector still reads healthy. Both detectors have
0/2{,}000 false positives (95\% upper bound 0.15\%). The profiler has 0/2{,}000
false negatives at full intensity; the peer-count-only detector has
2{,}000/2{,}000, since each attack preserves the clean-peer count.
\end{table}

Table~\ref{tab:p2p-detector} is a sensitivity analysis of the two detectors
against attack shapes that we define (Section~\ref{subsec:montecarlo}); it does
not measure robustness against arbitrary eclipse or Sybil strategies. Within
these four attacks the profiler flagged every full-intensity snapshot (0 misses
in 2{,}000) and every ramp (300 of 300 in each cell under three ramp shapes,
linear, front-loaded and slow; 95\% upper bound on the miss rate 1\%),
regardless of concurrent Bitcoin or DA degradation. It reads attacks of low
intensity as healthy: up to 0.25 for P4, 0.17 for P1 and 0.12 for P3. An
attacker who keeps every field within its bound is not detected by the six bounds
by construction (Section~\ref{subsec:sensors}). The peer-count-only detector
detected none of the 300 ramps in any cell.
